\begin{lemma}[Desigualdad de Van der Corput]
    \label{lem:vandercorput}
    Sean $H,N \in \mathbb{N}$ con $1 \leq H \leq N$ entonces,
    \begin{align}
        \left| \sum_{n=1}^N e^{2\pi i f(n)} \right|^2 \leq
        4 \, \frac{N}{H} \sum_{h=0}^{H-1} \left| \sum_{n=1}^{N-h} e^{2\pi i (f(n+h) - f(n))} \right|
    \end{align}
\end{lemma}

\begin{proof}
    Sea $u_n \in \mathbb{C}$ con $n = 1, \dots, N$. Definimos $u_n = 0$ para $n > N$ y 
    para $n \leq 0$. Entonces,
    \[
        H \sum_{n=1}^{N} u_n = \sum_{k=1}^{H} \sum_{n=1}^{N} u_n
    \]
    Sea $k$ fijo, consideramos el cambio de variable $p = n+k$ entonces $n = p-k$ y $p=1+k,\dots,N+k$
    \[
        \sum_{k=1}^{H} \sum_{n=1}^{N} u_n =
        \sum_{k=1}^{H} \sum_{p=k+1}^{N+k} u_{p-k}
    \]
    Por lo tanto, gracias a la extensión de ceros de los $u_n$,
    \[
        H \sum_{n=1}^{N} u_n =
        \sum_{k=1}^{H} \sum_{p=2}^{N+H} u_{p-k} =
        \sum_{p=2}^{N+H} \sum_{k=1}^{H} u_{p-k}
    \]
    
    Ahora elevamos al cuadrado ambos lados y aplicamos la desigualdad de Cauchy-Schwarz, \ref{ap:cauchy-schwarz},
    \[
        \left| H \sum_{n=1}^{N} u_n \right|^2 \leq
        (N+H-1) \sum_{p=2}^{N+H} \left| \sum_{k=1}^{H} u_{p-k} \right|^2
    \]
    
    Fijado $p$, vamos a desarrollar el sumatorio de dentro del término de la derecha,
    \[
        \left|  \sum_{k=1}^{H} u_{p-k} \right|^2 =
        \left(\sum_{k=1}^{H} u_{p-k} \right)\left( \sum_{s=1}^{H} \overline{u_{p-s}} \right) =
        \sum_{k=1}^{H} \sum_{s=1}^{H} u_{p-k} \overline{u_{p-s}}
    \]
    Este doble sumatorio se descompone en tres partes, la primera es la parte diagonal $k=s$,
    la segunda es la parte inferior $k < s$ y la tercera es la parte superior $k > s$,
    \[
        \sum_{k=1}^{H} \sum_{s=1}^{H} u_{p-k} \overline{u_{p-s}} =
        \sum_{k=1}^{H} |u_{p-k}|^2 \, +
        \sum_{1 \leq k < s \leq H} u_{p-k} \overline{u_{p-s}} \, +
        \sum_{1 \leq s < k \leq H} u_{p-k} \overline{u_{p-s}}
    \]
    Gracias a la simetría de los sumatorios de la parte inferior e superior,
    se pueden juntar ambos sumatorios y se obtiene,
    \[
        \sum_{k=1}^{H} \sum_{s=1}^{H} u_{p-k} \overline{u_{p-s}} =
        \sum_{k=1}^{H} |u_{p-k}|^2 \, +
        \sum_{1 \leq k < s \leq H} 2\operatorname{Re}(u_{p-k} \overline{u_{p-s}})
    \]

    De momento tenemos la siguiente cota,
    \begin{align*}
        H^2 \left| \sum_{n=1}^{N} u_n \right|^2 
        &\leq (N+H-1) \sum_{p=2}^{N+H} \sum_{k=1}^{H} |u_{p-k}|^2 +
        2(N+H-1) \sum_{p=2}^{N+H} \sum_{1 \leq k < s \leq H} \operatorname{Re}(u_{p-k} \overline{u_{p-s}}) \\
        &\leq (N+H-1) H \sum_{n=1}^{N} |u_{n}|^2 +
        2(N+H-1) \operatorname{Re}\left(\sum_{p=2}^{N+H} \sum_{1 \leq k < s \leq H} u_{p-k} \overline{u_{p-s}}\right)
    \end{align*}

    Estudiamos sumatorio del segundo término. Hacemos el cambio de variable $n = p-s$,
    notando que $u_n = 0$ salvo para $n=1,\dots,N$. Además, como la suma es finita podemos
    cambiar el orden de los sumatorios,
    \[
        \sum_{p=2}^{N+H} \sum_{1 \leq k < s \leq H} u_{p-k} \overline{u_{p-s}} =
        \sum_{1 \leq k < s \leq H} \sum_{n=1}^{N} u_{n+s-k} \overline{u_n}
    \]
    Sea $k$ fijo, hacemos el cambio de variable $h = s-k$. Como $k < s$ entonces 
    $s = k+1, \dots, H$ y $h = 1, \dots, H-k$. Obtenemos,
    \[
        \sum_{1 \leq k < s \leq H} \sum_{n=1}^{N} u_{n+s-k} \overline{u_n} =
        \sum_{k=1}^{H-1} \sum_{h=1}^{H-k} \sum_{n=1}^{N} u_{n+h} \overline{u_n}
    \]
    Sabemos que $1 \leq h \leq H-k$ entonces $2 \leq h+k \leq H$ y $1 \leq k \leq H-h$,
    \begin{align*}
        \sum_{k=1}^{H-1} \sum_{h=1}^{H-k} \sum_{n=1}^{N} u_{n+h} \overline{u_n}  
        &=\sum_{k=1}^{H-1} \sum_{h=1}^{H-1} \mathbf{1}_{[2,H]}(h+k) \left(\sum_{n=1}^{N} u_{n+h} \overline{u_n} \right) \\
        &=\sum_{h=1}^{H-1} \sum_{k=1}^{H-1} \mathbf{1}_{[2,H]}(h+k) \left(\sum_{n=1}^{N} u_{n+h} \overline{u_n} \right) \\
        &=\sum_{h=1}^{H-1} \sum_{k=1}^{H-h} \left(\sum_{n=1}^{N} u_{n+h} \overline{u_n} \right) \\
        &=\sum_{h=1}^{H-1} (H-h) \left(\sum_{n=1}^{N} u_{n+h} \overline{u_n} \right)
     \end{align*}

     La cota que teníamos la podemos simplificar como,
     \begin{align*}
        H^2 \left| \sum_{n=1}^{N} u_n \right|^2 
        &\leq (N+H-1) H \sum_{n=1}^{N} |u_{n}|^2 + 2(N+H-1) \operatorname{Re}\left(\sum_{h=1}^{H-1} (H-h) \left(\sum_{n=0}^{N} u_{n+h} \overline{u_n} \right)\right) \\        
        &\leq 2(N+H-1) \operatorname{Re}\left(\sum_{h=0}^{H-1} (H-h) \left(\sum_{n=0}^{N} u_{n+h} \overline{u_n} \right)\right) 
    \end{align*}

    Considerando las siguientes cotaciones $N+H-1 \leq 2N$, $H-h \leq H$ y 
    $\operatorname{Re}(z) \leq |z|$ para todo $z \in \mathbb{C}$ y teniendo en cuenta
    que $u_n=0$ salvo para $n=1,\dots,N$ obtenemos,
    \[
        H^2 \left| \sum_{n=1}^{N} u_n \right|^2 \leq
        4N H \sum_{h=0}^{H-1} \left| \sum_{n=1}^{N-h} u_{n+h} \overline{u_n} \right|
    \]

    Finalmente, dividiendo ambos lados por $H^2$ y aplicando $u_n = e^{2\pi i f(n)}$ 
    obtenemos la fórmula deseada,
    \[
        \left| \sum_{n=1}^N e^{2\pi i f(n)} \right|^2 \leq
        4 \, \frac{N}{H} \sum_{h=0}^{H-1} \left| \sum_{n=1}^{N-h} e^{2\pi i (f(n+h) - f(n))} \right|
    \]
\end{proof}